\documentclass[11pt, a4paper]{report}

\usepackage[utf8]{inputenc}
\usepackage[T1]{fontenc}
\usepackage[french]{babel}

% Appel des packages dans le dossier preambule
\def\packagepath{../../../../../preambule}   % path du package princial

\usepackage{\packagepath/page_layout} % <--- Nouveau
\usepackage{\packagepath/config_info}
\usepackage{\packagepath/cadres_cours}
\usepackage{\packagepath/zones_reponse}
\usepackage{\packagepath/config_ds}
\usepackage{\packagepath/config_maths}

% --- PILOTAGE ---
%\def\visibleornot{visible}    % visible or invisible
\ifdefined\visibleornot\else\def\visibleornot{visible}\fi


\def\documentpath{./Sources_Latex}
\graphicspath{{./Images/}}


\def\ptsexoA{0}



\begin{document}

\def\level{Premiere}  % Classe à droite
\def\course{NSI}
\def\chaptername{Constructions élémentaires} % Titre au centre
\def\docname{AP01~--~TD 1}        % Identifiant à gauche

\pagestyle{DS_LP}

\NomPrenom{}

\bigskip

%%=============================================================
\begin{exercice_cours}[Donner la valeur et le type de chacune des variables suivantes:]

    \begin{minipage}{0.48\linewidth}
        \begin{tabular}{p{2.6cm}cc}
    \hline
    Variable & Type & valeur\\
    \hline
    \verb|a = 4| & \reponseLigne[2cm]{int} & \reponseLigne[2cm]{\texttt{4}} \\
    \verb|b = 3.8| & \reponseLigne[2cm]{float} & \reponseLigne[2cm]{\texttt{3.8}} \\
    \verb|c = "hello"| & \reponseLigne[2cm]{str} & \reponseLigne[2cm]{\texttt{"hello"}} \\
    \verb|d = True| & \reponseLigne[2cm]{bool} & \reponseLigne[2cm]{\texttt{True}} \\
    \verb|e = a/2| & \reponseLigne[2cm]{float} & \reponseLigne[2cm]{\texttt{2.0}} \\
    \hline
        \end{tabular}
    \end{minipage}\hfill
    \begin{minipage}{0.48\linewidth}
        \begin{tabular}{p{2.6cm}cc}
    \hline
    Variable & Type & valeur\\
    \hline
    \verb|f = c[0]| & \reponseLigne[2cm]{str} & \reponseLigne[2cm]{\texttt{"h"}} \\
    \verb|g = len(c)| & \reponseLigne[2cm]{int} & \reponseLigne[2cm]{\texttt{5}} \\
    \verb|h = 15//2| & \reponseLigne[2cm]{int} & \reponseLigne[2cm]{\texttt{7}} \\
    \texttt{i = 15\%2} & \reponseLigne[2cm]{int} & \reponseLigne[2cm]{\texttt{1}} \\
    \verb|j = d and False| & \reponseLigne[2cm]{bool} & \reponseLigne[2cm]{\texttt{False}} \\
    
    \hline
        \end{tabular}
    \end{minipage}
\end{exercice_cours}
%%=============================================================

\bigskip

\begin{exercice_cours}[Donner ce qui est affiché après chaque instruction:]

    On considère les variables suivantes:

    \medskip

    \begin{minipage}{0.48\linewidth}
        \begin{lstlisting}[numbers=none]
    a = 4
    b = 2.5
    c = "hello"
    d = True
    \end{lstlisting}

    \end{minipage}\hfill{}
    \begin{minipage}{0.48\linewidth}
        \begin{lstlisting}[numbers=none]
    e = False
    f = '5'
    n = len(c)
    m = 15/3
    \end{lstlisting}

    \end{minipage}

    

    \begin{minipage}{0.48\linewidth}
        \begin{tabular}{p{3.5cm}c}
    \hline
    Variable &  valeur\\
    \hline
    \verb|str(a)| & \reponseLigne[4cm]{\texttt{'a'}} \\
    \verb|type(m)| & \reponseLigne[4cm]{\texttt{float}} \\
    \verb|c[1:3]| & \reponseLigne[4cm]{\texttt{'el'}} \\
    \verb|d and e| & \reponseLigne[4cm]{\texttt{False}} \\
    \verb|int(f)| & \reponseLigne[4cm]{\texttt{5}} \\
    \verb|'e' in c| & \reponseLigne[4cm]{\texttt{True}} \\
    \hline
        \end{tabular}
    \end{minipage}\hfill
    \begin{minipage}{0.48\linewidth}
        \begin{tabular}{p{3.5cm}c}
    \hline
    Variable &  valeur\\
    \hline
    \verb|c + f| & \reponseLigne[4cm]{\texttt{'hello5'}} \\
    \verb|n + a| & \reponseLigne[4cm]{\texttt{9}} \\
    \verb|a == 4| & \reponseLigne[4cm]{\texttt{True}} \\
    \verb|type(d)| & \reponseLigne[4cm]{\texttt{bool}} \\
    \verb|type(b) == str| & \reponseLigne[4cm]{\texttt{False}} \\
    \verb|'a' in c| & \reponseLigne[4cm]{\texttt{True}} \\
        \hline
        \end{tabular}
    \end{minipage}
\end{exercice_cours}

\medskip

\begin{exercice_cours}[On considère la variable \texttt{mot = 'informatique'}]
    Ecrire les codes python permettant d'obtenir les résultats suivants:
    
    \begin{minipage}{0.48\linewidth}
        \begin{tabular}{p{3.5cm}c}
    \hline
    résultat &  code python\\
    \hline
    \verb|'m'| & \reponseLigne[4cm]{\texttt{mot[5]}} \\
    \verb|'info'| & \reponseLigne[4cm]{\texttt{mot[0:4]}} \\
    \verb|'matique'| & \reponseLigne[4cm]{\texttt{mot[5:12]}} \\
    \verb|'tique'| & \reponseLigne[4cm]{\texttt{mot[7:12]}} \\
    \hline
        \end{tabular}
    \end{minipage}\hfill
    \begin{minipage}{0.48\linewidth}
        \begin{tabular}{p{3.5cm}c}
    \hline
    résultat &  code python\\
    \hline
    \verb|'o'| & \reponseLigne[4cm]{\texttt{mot[3]}} \\
    \verb|'nformatiqu'| & \reponseLigne[4cm]{\texttt{mot[1:-1]}} \\
    \verb|12| & \reponseLigne[4cm]{\texttt{len(mot)}} \\
    \verb|'e'| & \reponseLigne[4cm]{\texttt{mot[-1]}} \\
        \hline
        \end{tabular}
    \end{minipage}

\end{exercice_cours}

\end{document}